\chapter{集合与逻辑}
\label{cha:logic}

类型论，无论形式的还是非形式的，都是一套处理类型及其元素的规则。
然而，当用自然语言进行非形式的数学写作时，我们通常会使用一些熟悉的词汇，特别是``与''``或''等逻辑连接词，以及``对所有''``存在''等逻辑量词。
与集合论不同，类型论提供了不止一种方式来将这些英语短语理解为对类型的操作。
这种潜在的歧义需要加以消除，途径是制定局部或全局的约定、为非形式数学引入新的标注，或两者兼用。
这需要一定的适应过程，但好处在于：由于类型论允许对逻辑进行更精细的分析，我们能够比在集合论基础下更忠实地表达数学，减少``滥用语言''的情况。
本章将解释其中涉及的问题，并论证我们所做选择的合理性。

\section{集合与 \texorpdfstring{$n$}{n}-类型}
\label{sec:basics-sets}

\index{set|(defstyle}%

为了解释类型论的逻辑与集合论的逻辑之间的联系，在类型论中引入``集合''的概念是有益的。
一般而言，类型的行为类似于空间或高阶群胚，但其中有一个子类，其行为更接近于传统集合论系统中的集合。
从范畴论的角度看，我们可以考虑\emph{离散}群胚，它们由一组对象确定，且仅有恒同态射作为高阶态射；从拓扑学的角度看，我们可以考虑具有离散拓扑的空间。\index{discrete!space}
更一般地，我们可以考虑与这类结构\emph{等价}的群胚或空间；由于类型论中的一切都是在同伦意义下进行的，我们无法分辨它们的差异。

直觉上，如果一个类型不包含高阶同伦信息，我们就会预期它在这种意义下``是一个集合''：任何两条平行道路都是（在同伦意义下）相等的，所有更高维度的平行道路亦然。
幸运的是，由于同伦类型论中的一切构造都自动具有``函子性''/``连续性''，
\index{continuity of functions in type theory@``continuity'' of functions in type theory}%
\index{functoriality of functions in type theory@``functoriality'' of functions in type theory}%
实际上只需在最底层提出这一要求就足够了。

\begin{defn}\label{defn:set}
  一个类型 $A$ 是一个\define{集合（set）}，如果对于所有 $x,y:A$ 和所有 $p,q:x=y$，都有 $p=q$。
\end{defn}

更精确地说，命题 $\isset(A)$ 被定义为如下类型：
\[ \isset(A) \defeq \prd{x,y:A}{p,q:x=y} (p=q). \]

如 \cref{sec:types-vs-sets} 所述，同伦类型论中的集合不同于 ZF 集合论中的集合，因为不存在全局的``属于''谓词 $\in$。
它们更接近于结构数学和范畴论中使用的集合，其元素是``抽象的点''，我们通过函数和关系来赋予其结构。
这足以使它们成为大多数基于集合的数学的基础系统；我们将在 \cref{cha:set-math} 中看到一些例子。

哪些类型是集合？
在 \cref{cha:hlevels} 中我们将深入研究这个问题更一般的形式，但目前我们可以先看一些简单的例子。

\begin{eg}\label{eg:isset-unit}
  类型 \unit 是一个集合。
  因为根据 \cref{thm:path-unit}，对于任意 $x,y:\unit$，类型 $(x=y)$ 等价于 \unit。
  由于 \unit 的任意两个元素都相等，这意味着 $x=y$ 的任意两个元素也相等。
\end{eg}

\begin{eg}\label{eg:isset-empty}
  类型 $\emptyt$ 是一个集合，因为给定任意 $x,y:\emptyt$，根据 $\emptyt$ 的归纳原理，我们可以推导出任何结论。
\end{eg}

\begin{eg}\label{thm:nat-set}
  自然数类型 \nat 也是一个集合。
  这由 \cref{thm:path-nat} 可得，因为所有相等类型 $\id[\nat]xy$ 都等价于 \unit 或 $\emptyt$，而 \unit 或 $\emptyt$ 的任意两个实例都是相等的。
  我们将在 \cref{cha:hlevels} 中看到该事实的另一个证明。
\end{eg}

我们迄今考虑的大多数类型形成操作也都保持集合性质。

\begin{eg}\label{thm:isset-prod}
  如果 $A$ 和 $B$ 是集合，那么 $A\times B$ 也是集合。
  给定 $x,y:A\times B$ 和 $p,q:x=y$，根据 \cref{thm:path-prod} 有 $p= \pairpath(\projpath1(p),\projpath2(p))$ 和 $q= \pairpath(\projpath1(q),\projpath2(q))$。
  但由于 $A$ 是集合，可知 $\projpath1(p)=\projpath1(q)$；同理，因为 $B$ 是集合，可知 $\projpath2(p)=\projpath2(q)$。因此 $p=q$。

  类似地，如果 $A$ 是集合且 $B:A\to\type$ 满足每个 $B(x)$ 都是集合，那么 $\sm{x:A} B(x)$ 也是集合。
\end{eg}

\begin{eg}\label{thm:isset-forall}
  如果 $A$ 是\emph{任意}类型，且 $B:A\to \type$ 满足每个 $B(x)$ 都是集合，那么类型 $\prd{x:A} B(x)$ 也是一个集合。
  为此假设 $f,g:\prd{x:A} B(x)$ 且 $p,q:f=g$。
  根据函数外延性，我们有
  %
  \begin{equation*}
    p = {\funext (x \mapsto \happly(p,x))}
    \quad\text{and}\quad
    q = {\funext (x \mapsto \happly(q,x))}.
  \end{equation*}
  %
  但对任意 $x:A$，我们有
  %
  \begin{equation*}
   \happly(p,x):f(x)=g(x)
   \qquad\text{and}\qquad
   \happly(q,x):f(x)=g(x),
  \end{equation*}
  %
  由于 $B(x)$ 是集合，我们得到 $\happly(p,x) = \happly(q,x)$。
  现在再次使用函数外延性，依值函数 $(x \mapsto \happly(p,x))$ 与 $(x \mapsto \happly(q,x))$ 相等，因而（应用 $\apfunc{\funext}$）$p$ 与~$q$ 也相等。
\end{eg}

更多例子，参见 \cref{ex:isset-coprod,ex:isset-sigma}。关于同伦类型论中所有集合构成的子系统（范畴）的更系统研究，请参见~\cref{cha:set-math}。

\index{n-type@$n$-type|(}%

集合仅仅是所谓\emph{同伦 $n$-类型}阶梯上的第一级。
下一级由\emph{$1$-类型}构成，它们在范畴论中类似于 $1$-群胚。
集合（我们也可称之为\emph{$0$-类型}）的定义性质是其中没有非平凡的道路。
类似地，$1$-类型的定义性质是在道路之间没有非平凡的道路：

\begin{defn}\label{defn:1type}
  一个类型 $A$ 是\define{1-类型（1-type）}
  \indexdef{1-type}%
  当且仅当对所有 $x,y:A$ 以及 $p,q:x=y$ 和 $r,s:p=q$，都有 $r=s$。
\end{defn}

类似地，我们可以定义 $2$-类型、$3$-类型等等。
我们将在 \cref{cha:hlevels} 中归纳地定义 $n$-类型的一般概念，并研究不同 $n$ 值对应的 $n$-类型之间的关系。

然而，目前先了解两个事实是有益的。
第一，这些层级是向上封闭的：如果 $A$ 是 $n$-类型，那么 $A$ 也是 $(n+1)$-类型。
例如：
%%  will make precise the sense in which this
%% ``suffices for all higher levels'', but as an example, we observe that
%% it suffices for the next level up.

\begin{lem}\label{thm:isset-is1type}
  如果 $A$ 是集合（即 $\isset(A)$ 有实例），那么 $A$ 是 $1$-类型。
\end{lem}

\begin{proof}
  假设 $f:\isset(A)$；则对任意 $x,y:A$ 和 $p,q:x=y$，我们有 $f(x,y,p,q):p=q$。
  固定 $x$，$y$ 和 $p$，并定义 $g: \prd{q:x=y} (p=q)$ 为 $g(q) \defeq f (x,y,p,q)$。
  那么对任意 $r:q=q'$，我们有 $\apdfunc{g}(r) : \trans{r}{g(q)} = g(q')$。
  因此，由 \cref{cor:transport-path-prepost}，我们得到 $g(q) \ct r = g(q')$。

  特别地，假设给定如 \cref{defn:1type} 所示的 $x,y,p,q$ 和 $r,s:p=q$，并按上述定义 $g$。
  则 $g(p) \ct r = g(q)$ 且 $g(p) \ct s = g(q)$，故由消去律可得 $r=s$。
\end{proof}

第二，这种按层级对类型的划分并非退化，也就是说并非所有类型都是集合：

\begin{eg}\label{thm:type-is-not-a-set}
  \index{type!universe}%
  宇宙 \type 不是集合。
  为证明这一点，只需举出一个类型 $A$ 和一条道路 $p:A=A$，使得 $p$ 不等于 $\refl A$。
  取 $A=\bool$，并令 $f:A\to A$ 定义为 $f(\bfalse)\defeq \btrue$ 和 $f(\btrue)\defeq \bfalse$。
  那么对所有 $x$（通过简单的分情形讨论）都有 $f(f(x))=x$，因此 $f$ 是一个等价。
  于是，由泛等性，$f$ 给出了一条道路 $p:A=A$。

  如果 $p$ 等于 $\refl A$，那么（再次由泛等性）$f$ 将等于 $A$ 上的恒等函数。
  但这将意味着 $\bfalse=\btrue$，与 \cref{rmk:true-neq-false} 矛盾。
\end{eg}

在 \cref{cha:hits,cha:homotopy} 中我们将证明，对任意 $n$，都存在不是 $n$-类型的类型。

注意，$A$ 是 $1$-类型恰好意味着对任意 $x,y:A$，相等类型 $\id[A]xy$ 是集合。
（因此，\cref{thm:isset-is1type} 可以等价地理解为：集合的相等类型也是集合。）
这将是我们将在 \cref{cha:hlevels} 中给出的 $n$-类型递归定义的基础。

我们也可以将这一刻画从集合``向下''延伸。
也就是说，一个类型 $A$ 是集合当且仅当对任意 $x,y:A$，$\id[A]xy$ 中的任意两个元素都相等。
由于集合也就是 $0$-类型，因此自然称一个具有后一种性质（其任意两个元素相等）的类型为\emph{$(-1)$-类型}。
这类类型可视为\emph{狭义下的命题}，对它们的研究正是通常所谓的``逻辑''；这将占据本章剩余部分的内容。

\index{n-type@$n$-type|)}%
\index{set|)}%

\section{命题即类型？}
\label{subsec:pat?}

\index{proposition!as types|(}%
\index{logic!constructive vs classical|(}%
\index{anger|(}%
到目前为止，我们一直遵循着\cref{sec:pat}中描述的直接的``命题即类型''哲学。
按照这种哲学，像``存在 $x:A$ 使得 $P(x)$''这样的自然语言短语，被解释为对应的类型，例如 $\sm{x:A} P(x)$；
而一个命题的证明，则被视为判定某个特定元素居留于该类型。
然而，我们也已经看到，由这种解读所产生的``逻辑''，在某些方面会让熟悉经典数学的数学家感到陌生。
例如，在\cref{thm:ttac}中我们看到，形如
\index{axiom!of choice!type-theoretic}%
\begin{equation}\label{eq:english-ac}
  \parbox{\textwidth-2cm}{``If for all $x:X$ there exists an $a:A(x)$ such that $P(x,a)$, then there exists a function $g:\prd{x:X} A(x)$ such that for all $x:X$ we have $P(x,g(x))$'',}
\end{equation}
的陈述——它看上去很像经典的\index{mathematics!classical}\emph{选择公理（axiom of choice）}——在这种解读下总是真的。
这是命题即类型逻辑一个值得注意且通常很有用的特点，但它同时也表明了这种逻辑与逻辑的经典解释之间存在显著差异：
在经典解释下，选择公理并非逻辑真理，而是一个额外的``公理''。

另一方面，我们现在也可以证明，那些形如经典\emph{双重否定律（law of double negation）}和\emph{排中律（law of excluded middle）}的对应陈述，与泛等公理是不相容的。
\index{univalence axiom}%

\begin{thm}\label{thm:not-dneg}
  \index{double negation, law of}%
  并非对于所有 $A:\UU$ 都有 $\neg(\neg A) \to A$。
\end{thm}

\begin{proof}
  回顾 $\neg A \jdeq (A\to\emptyt)$。
  我们也将``并非……''读作算子 $\neg$。
  因此，为了证明这个陈述，只需假设给定某个$f:\prd{A:\UU} (\neg\neg A \to A)$，然后构造 \emptyt 的一个元素即可。

  接下来的证明思路是注意到，$f$ 和类型论中的任何函数一样，是``连续的''。
  \index{continuity of functions in type theory@``continuity'' of functions in type theory}%
  \index{functoriality of functions in type theory@``functoriality'' of functions in type theory}%
  由泛等公理可知，这意味着 $f$ 对于类型的等价是\emph{自然的}。
  由此，再加上一个无不动点的自等价\index{automorphism!fixed-point-free}，我们就能导出矛盾。

  令 $e:\eqv\bool\bool$ 为由 $e(\btrue)\defeq\bfalse$ 和 $e(\bfalse)\defeq\btrue$ 定义的等价，如\cref{thm:type-is-not-a-set}所示。
  令 $p:\bool=\bool$ 为由 $e$ 经泛等公理给出的对应道路，即 $p\defeq \ua(e)$。
  那么我们有 $f(\bool) : \neg\neg\bool \to\bool$，并且
  \[\apd f p : \transfib{A\mapsto (\neg\neg A \to A)}{p}{f(\bool)} = f(\bool).\]
  因此，对于任意 $u:\neg\neg\bool$，我们有
  \[\happly(\apd f p,u) : \transfib{A\mapsto (\neg\neg A \to A)}{p}{f(\bool)}(u) = f(\bool)(u).\]

  现在根据\eqref{eq:transport-arrow}，将 $f(\bool):\neg\neg\bool\to\bool$ 沿着 $p$ 在类型族${A\mapsto (\neg\neg A \to A)}$中传输，等于如下函数：
  它先将参数沿着 $\opp p$ 在类型族 $A\mapsto \neg\neg A$ 中传输，然后应用 $f(\bool)$，最后再将结果沿着 $p$ 在类型族 $A\mapsto A$ 中传输：
  \begin{narrowmultline*}
    \transfib{A\mapsto (\neg\neg A \to A)}{p}{f(\bool)}(u) =
    \narrowbreak
    \transfib{A\mapsto A}{p}{f(\bool) (\transfib{A\mapsto \neg\neg
        A}{\opp{p}}{u})}.
  \end{narrowmultline*}
  %
  然而，由函数外延性可知，任意两点 $u,v:\neg\neg\bool$ 都是相等的，因为对于任意 $x:\neg\bool$，我们有 $u(x):\emptyt$，从而可以推导出任何结论，特别是 $u(x)=v(x)$。
  因此，我们有$\transfib{A\mapsto \neg\neg A}{\opp{p}}{u} = u$，于是从 $\happly(\apd f p,u)$ 我们得到一个等式
  \[ \transfib{A\mapsto A}{p}{f(\bool)(u)} = f(\bool)(u).\]
  最后，如\cref{sec:compute-universe}所讨论的，沿着道路 $p\jdeq \ua(e)$ 在类型族 $A\mapsto A$ 中传输，等价于应用等价 $e$；因此我们有
  \begin{equation}
    e(f(\bool)(u)) = f(\bool)(u).\label{eq:fpaut}
  \end{equation}
  %
  然而，我们也可以证明
  \begin{equation}
    \prd{x:\bool} \neg(e(x)=x).\label{eq:fpfaut}
  \end{equation}
  这可以通过对 $x$ 进行分情形讨论得到：两种情形都能直接从 $e$ 的定义以及 $\bfalse\neq\btrue$（\cref{rmk:true-neq-false}）推出。
  因此，将\eqref{eq:fpfaut}应用于 $f(\bool)(u)$ 和\eqref{eq:fpaut}，我们就得到了 $\emptyt$ 的一个元素。
\end{proof}

\begin{rmk}
  \index{choice operator}%
  \indexsee{operator!choice}{choice operator}%
  特别地，这意味着不可能存在一个Hilbert 风格的``选择算子''，能够为每个非空类型选出一个元素。
  关键在于，任何这样的算子都不可能是\emph{自然的}，而在泛等公理下，所有作用于类型的函数对于等价都必须是自然的。
\end{rmk}

\begin{rmk}
  然而，对于任何 $A$，$\neg\neg\neg A \to \neg A$ 仍然成立；参见\cref{ex:neg-ldn}。
\end{rmk}

\begin{cor}\label{thm:not-lem}
  \index{excluded middle}%
  并非对于所有 $A:\UU$ 都有 $A+(\neg A)$。
\end{cor}

\begin{proof}
  假设我们有 $g:\prd{A:\UU} (A+(\neg A))$。
  我们将证明此时有 $\prd{A:\UU} (\neg\neg A \to A)$，从而可以应用 \cref{thm:not-dneg}。
  于是假设 $A:\UU$ 且 $u:\neg\neg A$；我们想要构造 $A$ 的一个元素。

  根据 $g(A):A+(\neg A)$，通过分情形讨论，我们可以假设要么 $g(A)\jdeq \inl(a)$ 对某个 $a:A$ 成立，要么 $g(A)\jdeq \inr(w)$ 对某个 $w:\neg A$ 成立。
  在第一种情形中，我们有 $a:A$；在第二种情形中，我们有 $u(w):\emptyt$，从而可以得出我们想要的任何结果（比如 $A$）。
  因此，在两种情形下我们都得到了 $A$ 的一个元素，如所需。
\end{proof}

因此，如果我们想假定泛等公理（我们当然想这样做），同时仍保留经典\index{mathematics!classical}推理的选项（这也是可取的），我们就不能使用未经修改的命题即类型原则来将\emph{所有}非形式的数学陈述解释到类型论中，因为那样排中律将为假。
然而，我们也不想完全抛弃命题即类型，因为它有许多优良性质（例如简单性、构造性以及可计算性）。
我们现在讨论一种对命题即类型的修改方案，以解决这些问题；在 \cref{subsec:when-trunc} 中，我们将回到何时使用哪种逻辑的问题。
\index{anger|)}%
\index{proposition!as types|)}%
\index{logic!constructive vs classical|)}%

\section{命题}
\label{subsec:hprops}

\index{logic!of mere propositions|(}%
\index{mere proposition|(defstyle}%
\indexsee{proposition!mere}{mere proposition}%
我们已经看到，命题即类型逻辑既有优点也有缺点。
两者有一个共同的原因：当类型被视作命题时，它们可以包含比单纯的真或假更多的信息，并且所有对它们的``逻辑''构造都必须尊重这些额外信息。
这表明，我们可以通过将注意力限制在那些\emph{不}包含任何超出真值的信息的类型上，并仅将这些类型视为逻辑命题，从而获得一种更传统的逻辑。

这样的类型 $A$ 如果有实例则视为``真''，如果从其实例可导出矛盾（即如果 $\neg A \jdeq (A\to\emptyt)$ 是有实例的）则视为``假''。
\index{inhabited type}%
为了获得更传统的逻辑，我们希望避免将这样的类型当作逻辑命题：给出其一个元素会提供比仅知道该类型有实例更多的信息。
例如，如果给定 \bool 的一个元素，那么我们得到的信息就比仅仅知道 \bool 包含某个元素要多。
事实上，我们正好多得到了\emph{一比特}\index{bit}的信息：我们知道所给的是 \bool 的\emph{哪一个}元素。
相比之下，如果给定 \unit 的一个元素，那么我们得到的信息不比仅知道 \unit 包含一个元素更多，因为 \unit 的任意两个元素都彼此相等。
这启发我们做出如下定义。

\begin{defn}\label{defn:isprop}
  一个类型 $P$ 称为\define{命题（mere proposition）}，如果对于所有 $x,y:P$ 都有 $x=y$。
\end{defn}

注意，由于我们仍然\emph{在}类型论中做数学，所以这也是一个\emph{在}类型论中的定义，这意味着它本身是一个类型——或者更准确地说，是一个类型族。
具体来说，对于任意 $P:\type$，类型 $\isprop(P)$ 定义为
\[ \isprop(P) \defeq \prd{x,y:P} (x=y). \]
因此，断言``$P$ 是一个命题''就是要给出 $\isprop(P)$ 的一个实例，即一个依值函数，它通过道路将 $P$ 的任意两个元素连接起来。
这个函数的连续性/自然性意味着，不仅 $P$ 的任意两个元素相等，而且 $P$ 也不包含任何高阶同伦。

\begin{lem}\label{thm:inhabprop-eqvunit}
  如果 $P$ 是一个命题且 $x_0:P$，则有 $\eqv P \unit$。
\end{lem}

\begin{proof}
  定义 $f:P\to\unit$ 为 $f(x)\defeq \ttt$，定义 $g:\unit\to P$ 为 $g(u)\defeq x_0$。
  结论由下一个引理以及 \unit 是一个命题（根据 \cref{thm:path-unit}）这一观察得出。
\end{proof}

\begin{lem}\label{lem:equiv-iff-hprop}
  如果 $P$ 和 $Q$ 都是命题，且 $P\to Q$ 与 $Q\to P$ 成立，那么 $\eqv P Q$。
\end{lem}

\begin{proof}
  假设给定 $f:P\to Q$ 和 $g:Q\to P$。
  那么对于任意 $x:P$，由于 $P$ 是命题，我们有 $g(f(x))=x$。
  类似地，对于任意 $y:Q$，由于 $Q$ 是命题，我们有 $f(g(y))=y$；因此 $f$ 和 $g$ 互为拟逆。
\end{proof}

也就是说，正如在 \cref{sec:pat} 中所承诺的，如果两个命题是逻辑等价的，那么它们就是等价的。

在同伦论中，一个与 \unit 同伦等价的空间被称为\emph{可缩的}。
因此，任何有实例的命题都是可缩的（另见\cref{sec:contractibility}）。
另一方面，空类型 \emptyt 也是一个（平凡地）命题。
至少在经典\index{mathematics!classical}数学中，只有这两种可能性。

命题也称为\emph{子终对象}（从范畴论角度）、\emph{子单点集}（从集合论角度），或\emph{h-命题}。
\indexsee{object!subterminal}{mere proposition}%
\indexsee{subterminal object}{mere proposition}%
\indexsee{subsingleton}{mere proposition}%
\indexsee{h-proposition}{mere proposition}
\cref{sec:basics-sets}中的讨论提示我们也可以称其为\emph{$(-1)$-类型}；我们将在\cref{cha:hlevels}回到这一点。
形容词``mere''强调的是：尽管任何类型都可以被视作一个命题（我们通过给出它的一个实例来做到这一点），但一个仅仅是命题的类型，不能被有用地视作比命题\emph{更多}的东西：在其为真的见证中不包含任何额外信息。

注意，一个类型 $A$ 是集合当且仅当对任意 $x,y:A$，相等类型 $\id[A]xy$ 都是一个命题。
另一方面，通过复制并简化\cref{thm:isset-is1type}的证明，我们得到：

\begin{lem}\label{thm:prop-set}
  每个命题都是集合。
\end{lem}

\begin{proof}
  假设 $f:\isprop(A)$；于是对任意 $x,y:A$ 有 $f(x,y):x=y$。固定 $x:A$ 并定义$g(y)\defeq f(x,y)$。那么对任意 $y,z:A$ 和 $p:y=z$，我们有$\apd
  g p : \trans{p}{g(y)}={g(z)}$。因此由\cref{cor:transport-path-prepost}，我们得到$g(y)\ct p = g(z)$，这就是说$p=\opp{g(y)}\ct g(z)$。于是，对任意 $p,q:x=y$，我们有$p = \opp{g(x)}\ct g(y) = q$。
\end{proof}

特别地，这意味着：

\begin{lem}\label{thm:isprop-isprop}\label{thm:isprop-isset}
  对任意类型 $A$，类型 $\isprop(A)$ 和 $\isset(A)$ 都是命题。
\end{lem}

\begin{proof}
  假设 $f,g:\isprop(A)$。由函数外延性，要证明 $f=g$，只需证明对任意 $x,y:A$ 有$f(x,y)=g(x,y)$。而 $f(x,y)$ 和 $g(x,y)$ 都是 $A$ 中的道路，它们之所以相等，是因为由 $f$ 或 $g$ 可知 $A$ 是一个命题，从而由\cref{thm:prop-set}可知它是一个集合。
  类似地，假设 $f,g:\isset(A)$，即对任意 $a,b:A$ 和 $p,q:a=b$，我们有$f(a,b,p,q):p=q$和$g(a,b,p,q):p=q$。但由于 $A$ 是一个集合（由 $f$ 或 $g$ 可知），故而是一个 $1$-类型，由此可得$f(a,b,p,q)=g(a,b,p,q)$；于是由函数外延性有 $f=g$。
\end{proof}

到目前为止我们还见过另一个例子：\cref{sec:basics-equivalences}中的条件~\ref{item:be3}断言，对任意函数 $f$，类型 $\isequiv (f)$ 应该是一个命题。

\index{logic!of mere propositions|)}%
\index{mere proposition|)}%

\section{经典逻辑与直觉主义逻辑}
\label{sec:intuitionism}

\index{logic!constructive vs classical|(}%
\index{denial|(}%
有了命题的概念，我们现在可以给出同伦类型论中\define{排中律（law of excluded middle）}
\indexdef{excluded middle}%
\indexsee{axiom!excluded middle}{excluded middle}%
\indexsee{law!of excluded middle}{excluded middle}%
的正确定义：
\begin{equation}
  \label{eq:lem}
  \LEM{}\;\defeq\;
  \prd{A:\UU} \Big(\isprop(A) \to (A + \neg A)\Big).
\end{equation}
类似地，\define{双重否定律（law of double negation）}
\indexdef{double negation, law of}%
\indexdef{axiom!double negation}%
\indexdef{law!of double negation}%
是
\begin{equation}
  \label{eq:ldn}
  % \mathsf{DN}\;\defeq\;
  \prd{A:\UU} \Big(\isprop(A) \to (\neg\neg A \to A)\Big).
\end{equation}
二者也很容易看出是等价的——参见\cref{ex:lem-ldn}——因此从现在起我们通常只提 \LEM{}。

\LEM{} 的这一表述避免了\cref{thm:not-dneg,thm:not-lem}中的``悖论''，因为 \bool 不是一个命题。
为了将其与更一般的``命题即类型''表述区分开来，我们将后者重命名为：
\symlabel{lem-infty}
\begin{equation*}
  \LEM\infty \defeq \prd{A:\UU} (A + \neg A).
\end{equation*}
为强调起见，正确的版本~\eqref{eq:lem}
可以记作 $\LEM{-1}$；
另见\cref{ex:lemnm}。
虽然 $\LEM{}$
不是\cref{cha:typetheory}中所描述的基础类型论的推论，但可以一致地将其假定为一公理（不像它的 $\infty$ 版本）。
例如，我们将在\cref{sec:wellorderings}中假设它。

然而，不使用 \LEM{} 也能走很远，这一点可能会令人惊讶。
通常，只需对定义或定理稍加改写，就能避免诉诸排中律。
虽然这有时需要一点时间适应，但这番周折往往是值得的，能带来更优雅、更一般的证明。
我们在引论中已经讨论了这样做的一些益处。

例如，在经典\index{mathematics!classical}数学中，双重否定常被不必要地使用。
一个非常简单的例子是通常的假设：集合 $A$ 是“非空的”，其字面意思是“$A$ \emph{不}包含\emph{任何}元素”这一情形不成立。
而几乎总是，其真正含义是 $A$ \emph{确实}包含至少一个元素这一正面断言；去掉这层双重否定，我们就能让陈述减少对 \LEM{} 的依赖。
回想一下，当我们把 $A$ 本身作为命题来断言时（即构造出 $A$ 的一个元素，通常不命名），我们就说类型 $A$ 是\emph{有实例的（inhabited）}
\index{inhabited type}%
。
因此，在将经典证明转译为构造逻辑时，我们常常将“非空的”替换为“有实例的”（尽管有时必须替换为“仅知有实例的”；参见\cref{subsec:prop-trunc}）。

类似地，在经典数学中，不必要的反证法也并不罕见。
\index{proof!by contradiction}%
当然，经典反证法的形式依赖双重否定律：我们先假设 $\neg A$ 并推导出矛盾，从而得到 $\neg \neg A$，再由双重否定律得到 $A$。
然而，从 $\neg A$ 推导出矛盾的过程，常常稍作改写即可给出 $A$ 的一个直接证明，从而无需使用 \LEM{}。

同样需要注意的是，如果目标是证明一个\emph{否定}\index{negation}，那么“反证法”并不涉及 \LEM{}。
事实上，$\neg A$ 根据定义就是类型 $A\to\emptyt$，所以证明 $\neg A$ 无非是在假设 $A$ 的前提下证明矛盾（$\emptyt$）。
类似地，双重否定律对否定命题确实成立：$\neg\neg\neg A \to \neg A$。
通过练习，我们会逐渐学会更仔细地区分否定命题与非否定命题，并注意到何时用到了 \LEM{}、何时没有。

因此，与表面看来相反，“构造性地”做数学通常并不意味着放弃重要定理，而是要找到表述定义的最佳方式，使重要定理能够构造性地证明。
也就是说，我们在初步研究一个主题时可以自由使用 \LEM{}，但一旦对该主题有了更深入的理解，就有望改进其定义和证明，从而不再依赖那条公理。
% For instance, the theory of ordinal numbers, which classically makes heavy use of \LEM{}, works quite well constructively once we choose the correct definition of ``ordinal''; see \cref{sec:ordinals}.
这一观察在同伦类型论中尤为显著：泛等性和高阶归纳类型这些强大工具使我们能够构造性地应对许多传统上需要经典\index{mathematics!classical}推理的问题。
我们将在\cref{part:mathematics}中看到若干这样的例子。
% For instance, none of the ``synthetic'' homotopy theory we will develop in \cref{cha:homotopy} requires \LEM{} --- despite the fact that classical homotopy theory (formulated using topological spaces or simplicial sets) makes heavy use of them (as well as the axiom of choice).

还值得提及的是，即便在构造性数学中，排中律对\emph{某些}命题也可以成立。
这类命题传统上称为\emph{可判定的（decidable）}。

\begin{defn}\label{defn:decidable-equality}
  \mbox{}
  \begin{enumerate}
  \item 若 $A+\neg A$，则称类型 $A$ 为\define{可判定的（decidable）}
    \indexdef{decidable!type}%
    \indexdef{type!decidable}%
    。
  \item 类似地，若 \narrowequation{\prd{a:A} (B(a)+\neg B(a))}， 则称类型族 $B:A\to \type$ 是\define{可判定的（decidable）}
    \indexdef{decidable!type family}%
    \indexdef{type!family of!decidable}%
    。 \label{item:decidable-equality2}
  \item 特别地，若 \narrowequation{\prd{a,b:A} ((a=b) + \neg(a=b))}， 则称 $A$ 具有\define{可判定相等性（decidable equality）}
    \indexdef{decidable!equality}%
    \indexsee{equality!decidable}{decidable equality}%
    。
  \end{enumerate}
\end{defn}

因此，$\LEM{}$ 恰好断言了所有命题都是可判定的，从而所有命题族也是可判定的。
特别地，$\LEM{}$ 蕴含所有集合（在\cref{sec:basics-sets}的意义下）都具有可判定相等性。
这种意义上的可判定相等性是非常强的性质；参见\cref{thm:hedberg}。

\index{denial|)}%
\index{logic!constructive vs classical|)}%

\section{子集与命题大小重整}
\label{subsec:prop-subsets}

\index{mere proposition|(}%

作为命题有用性的另一个例子，我们来讨论子集（以及更一般的子类型）。
假设 $P:A\to\type$ 是一个类型族，且每个类型 $P(x)$ 都被视为一个命题。
那么 $P$ 本身就是 $A$ 上的一个\emph{谓词（predicate）}，或者说 $A$ 的元素的一个\emph{性质（property）}。

在集合论中，给定集合$A$上的一个谓词$P$，我们便可以构造子集$\setof{x\in A | P(x)}$。
如\cref{sec:pat}中所简要提及的，在类型论中与之最直接对应的类比便是$\Sigma$-类型 $\sm{x:A} P(x)$。
当然，$\sm{x:A} P(x)$的一个实例是一个对子$(x,p)$，其中$x:A$且$p$是$P(x)$的一个证明。
然而，对于一般的$P$，一个元素$a:A$可能会产生$\sm{x:A} P(x)$中多个不同的元素——如果命题$P(a)$有多个不同的证明。
这就有悖于我们通常对子集的直觉。
但倘若$P$是一个\emph{纯}命题（mere proposition），那么这种情况便不会发生。

\begin{lem}\label{thm:path-subset}
  设$P:A\to\type$是一个类型族，且对任意$x:A$，$P(x)$均为一个命题。
  若$u,v:\sm{x:A} P(x)$满足$\proj1(u) = \proj1(v)$，则有$u=v$。
\end{lem}

\begin{proof}
  假设$p:\proj1(u) = \proj1(v)$。
  根据\cref{thm:path-sigma}，为证明$u=v$，只需证$\trans{p}{\proj2(u)} = \proj2(v)$。
  然而，$\trans{p}{\proj2(u)}$与$\proj2(v)$都是$P(\proj1(v))$的元素，而$P(\proj1(v))$是命题；因此它们是相等的。
\end{proof}

举例来说，回忆我们在\cref{sec:basics-equivalences}中定义了
\[(\eqv A B) \;\defeq\; \sm{f:A\to B} \isequiv (f),\]
其中每个类型$\isequiv (f)$都被假定为一个命题。
由此可知，若两个等价的底层函数相等，则它们作为等价而言也是相等的。

\label{defn:setof}%
此后，若$P:A\to \type$是一个由命题组成的族（即每个$P(x)$都是命题），我们可以将
%
\begin{equation}
  \label{eq:subset}
  \setof{x:A | P(x)}
\end{equation}
%
作为$\sm{x:A} P(x)$的另一种记法。
（在技术上，没有理由不能将此记法也用于一般的$P$，但这种用法可能会因带来意外的隐含联想而引起混淆。）
若$A$是一个集合，我们称\eqref{eq:subset}为$A$的一个\define{子集}
\indexdef{subset}%
\indexdef{type!subset}%
；对于一般的$A$，我们或可称其为\define{子类型}。
\indexdef{subtype}%
我们也可以直接将$P$本身称作$A$的一个\emph{子集}或\emph{子类型}；这实际上更为正确，因为孤立地看类型~\eqref{eq:subset}并不记得它与$A$的关系。

\symlabel{membership}
对于这样的$P$和$a:A$，我们可以记作$a\in P$或$a\in \setof{x:A | P(x)}$以指代命题$P(a)$。
若它成立，则我们称$a$是$P$的一个\define{元素}。
\symlabel{subset}
类似地，若$\setof{x:A | Q(x)}$是$A$的另一个子集，则说$P$\define{被包含}
\indexdef{containment!of subsets}%
\indexdef{inclusion!of subsets}%
于$Q$中，并记作$P\subseteq Q$，若有$\prd{x:A}(P(x)\rightarrow Q(x))$。

作为子类型的更多例子，我们可以定义某个宇宙\UU{}中所有集合与所有命题构成的“子宇宙”：
\symlabel{setU}\symlabel{propU}
\begin{align*}
  \setU &\defeq \setof{A:\UU | \isset(A) },\\
  \propU &\defeq \setof{A:\UU | \isprop(A) }.
\end{align*}
$\setU$的一个元素是一个类型$A:\UU$连同证据$s:\isset(A)$，$\propU$亦同理。
\cref{thm:path-subset}蕴含着$\id[\setU]{(A,s)}{(B,t)}$等价于$\id[\UU]AB$（因而也等价于$\eqv AB$）。
因此，我们经常会滥用记号，简写$A:\setU$而非$(A,s):\setU$。
若无需指定所讨论的具体宇宙，我们也可以省略下标\UU。

回忆一下，对于任意两个宇宙$\UU_i$和$\UU_{i+1}$，若$A:\UU_i$则亦有$A:\UU_{i+1}$。
因此，对于任意$(A,s):\set_{\UU_i}$，我们也有$(A,s):\set_{\UU_{i+1}}$，对$\prop_{\UU_i}$亦类似，从而得到自然的映射
\begin{align}
  \set_{\UU_i} &\to \set_{\UU_{i+1}}\label{eq:set-up},\\
  \prop_{\UU_i} &\to \prop_{\UU_{i+1}}.\label{eq:prop-up}
\end{align}
映射~\eqref{eq:set-up}不可能是一个等价，否则我们将重现康托尔集合论中为人熟知的自指悖论\index{paradox}。
然而，尽管在目前所给出的类型论中，~\eqref{eq:prop-up}并非自动成为一个等价，但假定其为一个等价是一致的。
也就是说，我们可以考虑向类型论中添加如下公理。

\begin{axiom}[命题大小重整]
  \indexsee{axiom!propositional resizing}{propositional resizing}%
  \indexdef{propositional!resizing}%
  \indexsee{resizing!propositional}{propositional resizing}%
  映射$\prop_{\UU_i} \to \prop_{\UU_{i+1}}$是一个等价。
\end{axiom}

我们将此公理称为\define{命题大小重整}，
因为它意味着宇宙$\UU_{i+1}$中的任何命题均可被“重整大小”为较小宇宙$\UU_i$中与之等价的一个命题。
若$\UU_{i+1}$满足\LEM{}，则此公理自动成立（参见\cref{ex:lem-impred}）。
我们一般不假定此公理成立，不过在有些地方会将其用作明确的假设。
它是针对命题的一种\emph{非直谓性}形式，而避免使用它，则称类型论保持\emph{直谓的}。
\indexsee{impredicativity!for mere propositions}{propositional resizing}%
\indexdef{mathematics!predicative}%

在实践中，我们最常需要的是一个略有不同的陈述：所考虑的宇宙 \UU 包含一个能``囊括所有命题''的类型。
换言之，我们需要一个类型 $\Omega:\UU$ 以及一族以 $\Omega$ 为指标的命题，使得该族在等价的意义下包含所有命题。
如果 \UU 不是最小的宇宙 $\UU_0$，那么上述命题大小重整即可推出此陈述，因为此时我们可以定义 $\Omega\defeq \prop_{\UU_0}$。

非直谓性的一个用途是定义幂集。
将集合 $A$ 的\define{幂集（power set）}\indexdef{power set}定义为 $A\to\propU$ 是很自然的；但在缺乏非直谓性的情况下，这个定义（哪怕仅就等价而言）也依赖于宇宙 \UU 的选取。
然而，有了命题大小重整，我们就可以将幂集定义为
\symlabel{powerset}%
\[ \power A \defeq (A\to\Omega),\]
这样它就不依赖于 $\UU$ 了。
另见 \cref{subsec:piw}。

\section{命题的逻辑}
\label{subsec:logic-hprop}

\index{logic!of mere propositions|(}%
我们曾在 \cref{sec:types-vs-sets} 提到，与仅有一个基本概念（类型）的类型论不同，集合论基础有两个基本概念：集合与命题。
因此，一位经典\index{mathematics!classical}数学家习惯于分开处理这两种对象。

在类型论中也可以恢复类似的二分法：由那些是\emph{命题}的类型（及类型族）来扮演集合论中命题的角色。
在许多情况下，只需将相应的类型形成子限制到命题上，就可以在这种逻辑中表示出逻辑连接词和量词。
当然，这要求我们知道所讨论的类型形成子能保持命题性。

\begin{eg}
  如果 $P$ 和 $Q$ 是命题，那么 $P\times Q$ 也是命题。
  这很容易证明，只需利用积类型中道路的刻画，类似于 \cref{thm:isset-prod} 但更简单。
  因此，连接词``且''保持命题性。
\end{eg}

\begin{eg}\label{thm:isprop-forall}
  如果 $A$ 是任意类型，且 $P:A\to \type$ 满足对一切 $x:A$，类型 $P(x)$ 都是命题，那么 $\prd{x:A} P(x)$ 也是命题。
  证明思路类似于 \cref{thm:isset-forall} 但更简单：给定 $f,g:\prd{x:A} P(x)$，对任意 $x:A$，由于 $P(x)$ 是命题，我们有 $f(x)=g(x)$。
  然后根据函数外延性，即得 $f=g$。

  特别地，如果 $P$ 是命题，那么无论 $A$ 是什么，$A\to P$ 也是命题。
  更进一步，由于 \emptyt 是命题，所以 $\neg A \jdeq (A\to\emptyt)$ 也是命题。
  \index{quantifier!universal}%
  因此，连接词``蕴含''和``非''保持命题性，量词``对所有''也是如此。
\end{eg}

另一方面，并非所有类型形成子都能保持命题。
即使 $P$ 和 $Q$ 是命题， $P+Q$ 在一般情况下也未必是。
例如， \unit 是一个命题，但 $\bool=\unit+\unit$ 不是。
从逻辑上讲， $P+Q$ 是一种``纯粹构造性''的``或''：它的一个见证会额外包含\emph{哪个}析取支为真的信息。
有时这非常有用，但如果我们想要一种更经典的、能保持命题的``或''，就需要一种方法来``截断''这个类型，通过遗忘这些额外信息将其变为一个命题。

\index{quantifier!existential}%
$\Sigma$-类型 $\sm{x:A} P(x)$（其中 $A$ 是任意类型）也存在同样的问题。
这是``存在一个 $x:A$ 使得 $P(x)$''的纯粹构造性解释，它记住了见证 $x$，因此即使每个类型 $P(x)$ 都是命题，该类型本身通常也不是。
（回想一下，我们在 \cref{subsec:prop-subsets} 中曾指出， $\sm{x:A} P(x)$ 也可以看作``满足 $P(x)$ 的那些 $x:A$ 构成的子集''。）

\section{命题截断}
\label{subsec:prop-trunc}

\index{truncation!propositional|(defstyle}%
\indexsee{type!squash}{truncation, propositional}%
\indexsee{squash type}{truncation, propositional}%
\indexsee{bracket type}{truncation, propositional}%
\indexsee{type!bracket}{truncation, propositional}%
\emph{命题截断（propositional truncation）}，也称为\emph{$(-1)$-截断}、\emph{括号类型（bracket type）}或\emph{挤压类型（squash type）}，是一个额外的类型形成子，它能将一个类型``挤压''或``截断''为一个命题，从而忘掉该类型中居民所包含的、除其存在以外的所有信息。

更准确地说，对于任意类型 $A$，都存在一个类型 $\brck{A}$。
它有两个构造子：

\begin{itemize}
\item 对于任意 $a:A$，我们有 $\bproj a : \brck A$。
\item 对于任意 $x,y:\brck A$，我们有 $x=y$。
\end{itemize}

第一个构造子意味着，如果 $A$ 是有实例的，那么 $\brck A$ 也是有实例的。
第二个构造子确保 $\brck A$ 是一个命题；通常我们不给这个事实的见证命名。

\index{recursion principle!for truncation}%
$\brck A$ 的递归原理如下：

\begin{itemize}
\item 如果 $B$ 是一个命题，并且我们有 $f:A\to B$，那么就存在一个诱导出的 $g:\brck A \to B$，使得对所有 $a:A$，都有 $g(\bproj a) \jdeq f(a)$。
\end{itemize}

换句话说，任何（从 $A$ 的有实例性）能推出的命题，也都能从 $\brck A$ 推出。
因此，作为命题的 $\brck A$，所包含的信息并不比 $A$ 的有实例性更多。
（$\brck A$ 也有一个归纳原理，但它不是特别有用；参见 \cref{ex:prop-trunc-ind}。）

在 \cref{ex:lem-brck,ex:impred-brck,sec:hittruncations} 中，我们将描述如何用更一般的概念来构造 $\brck{A}$。
目前，我们只是将其作为 \cref{cha:typetheory} 中那些规则之外的一条附加规则来使用。

有了命题截断，我们就可以将``命题的逻辑''扩展到涵盖析取与存在量词。
具体来说，$\brck{A+B}$ 是``$A$ 或 $B$''的命题版本，它并不``记住''哪个析取支为真的信息。

截断的递归原理意味着，\emph{当我们试图证明一个命题时}，仍然可以在 $\brck{A+B}$ 上进行分情形讨论。
也就是说，假设我们有 $u:\brck{A+B}$，并且我们试图证明一个命题 $Q$。
换句话说，我们试图定义 $\brck{A+B} \to Q$ 的一个元素。
由于 $Q$ 是一个命题，根据命题截断的递归原理，我们只需构造一个函数 $A+B\to Q$ 即可。
于是现在我们可以对 $A+B$ 进行分情形讨论了。

类似地，对于类型族 $P:A\to\type$，我们可以考虑 $\brck{\sm{x:A} P(x)}$，它是``存在 $x:A$ 使得 $P(x)$''的命题版本。
与析取的情形一样，通过结合截断与 $\Sigma$ 类型的归纳原理，如果我们有一个类型为 $\brck{\sm{x:A} P(x)}$ 的假设，那么\emph{当试图证明一个命题时}，我们可以引入新的假设 $x:A$ 和 $y:P(x)$。
换句话说，如果我们知道存在某个 $x:A$ 使得 $P(x)$，但我们手头并没有一个具体的这样的 $x$，那么我们仍然可以自由地使用这样的 $x$，只要我们不是在试图构造任何可能依赖于 $x$ 的具体取值的对象。
要求陪域是一个命题就表达了结果对见证的独立性，因为这样一个类型的所有可能的实例都必须相等。

在\cref{cha:real-numbers,cha:set-math}所讨论的集合层面的数学中，我们主要处理集合与命题，此时使用传统的逻辑记号来专门指称``命题截断逻辑''会很方便。

\begin{defn} \label{defn:logical-notation}
  我们通过截断来定义\define{传统逻辑记号}
  \indexdef{implication}%
  \indexdef{traditional logical notation}%
  \indexdef{logical notation, traditional}%
  \index{quantifier}%
  \indexsee{existential quantifier}{quantifier, existential}%
  \index{quantifier!existential}%
  \indexsee{universal!quantifier}{quantifier, universal}%
  \index{quantifier!universal}%
  \indexdef{conjunction}%
  \indexdef{disjunction}%
  \indexdef{true}%
  \indexdef{false}%
  ，其中 $P$ 和 $Q$ 表示命题（或命题族）：
  {\allowdisplaybreaks
  \begin{align*}
    \top            &\ \defeq \ \unit \\
    \bot            &\ \defeq \ \emptyt \\
    P \land Q       &\ \defeq \ P \times Q \\
    P \Rightarrow Q &\ \defeq \ P \to Q \\
    P \Leftrightarrow Q &\ \defeq \ P = Q \\
    \neg P          &\ \defeq \ P \to \emptyt \\
    P \lor Q        &\ \defeq \ \brck{P + Q} \\
    \fall{x : A} P(x) &\ \defeq \ \prd{x : A} P(x) \\
    \exis{x : A} P(x) &\ \defeq \ \Brck{\sm{x : A} P(x)}
  \end{align*}}
\end{defn}

符号 $\land$ 和 $\lor$ 在同伦论中也用于表示带点空间的缩积与楔，我们将在\cref{cha:hits}引入它们。严格来说这造成了潜在的冲突，但通常不会引起混淆。

类似地，当如\cref{subsec:prop-subsets}那样讨论子集时，我们也会使用传统的记号来表示交、并、补：
\indexdef{intersection!of subsets}%
\symlabel{intersection}%
\indexdef{union!of subsets}%
\symlabel{union}%
\indexdef{complement, of a subset}%
\symlabel{complement}%
\begin{align*}
  \setof{x:A | P(x)} \cap \setof{x:A | Q(x)}
  &\defeq \setof{x:A | P(x) \land Q(x)},\\
  \setof{x:A | P(x)} \cup \setof{x:A | Q(x)}
  &\defeq \setof{x:A | P(x) \lor Q(x)},\\
  A \setminus \setof{x:A | P(x)}
  &\defeq \setof{x:A | \neg P(x)}.
\end{align*}
当然，在没有排中律的情况下，后者并非通常意义上的``补集''：对于 $A$ 的任意子集 $B$，我们未必有 $B \cup (A\setminus B) = A$。

\index{truncation!propositional|)}%
\index{mere proposition|)}%
\index{logic!of mere propositions|)}%

\section{选择公理}
\label{sec:axiom-choice}

\index{axiom!of choice|(defstyle}%
\index{denial|(}%
现在我们可以恰当表述同伦类型论中的选择公理了。
假设给定一个类型 $X$ 以及类型族
%
\begin{equation*}
  A:X\to\type
  \qquad\text{and}\qquad
  P:\prd{x:X} A(x)\to\type,
\end{equation*}
%
并且满足以下条件：

\begin{itemize}
\item $X$ 是一个集合；
\item 对每个 $x:X$，$A(x)$ 是一个集合；
\item 对每个 $x:X$ 和 $a:A(x)$，$P(x,a)$ 是一个命题。
\end{itemize}

\define{选择公理（axiom of choice）}
$\choice{}$ 断言，在如下假设之下，
\begin{equation}\label{eq:ac}
  \Parens{\prd{x:X} \Brck{\sm{a:A(x)} P(x,a)}}
  \to
  \Brck{\sm{g:\prd{x:X} A(x)} \prd{x:X} P(x,g(x))}.
\end{equation}
当然，这是~\eqref{eq:english-ac} 的直接翻译，其中我们把``存在 $x:A$ 使得 $B(x)$''读作 $\brck{\sm{x:A}B(x)}$，所以也可以用熟悉的逻辑记号将这一陈述写作
\begin{narrowmultline*}
  \textstyle
  \Big(\fall{x:X}\exis{a:A(x)} P(x,a)\Big)
  \Rightarrow \narrowbreak
  \Big(\exis{g : \prd{x:X} A(x)} \fall{x : X} P(x,g(x))\Big).
\end{narrowmultline*}
%
特别要注意，命题截断在这里出现了两次。
定义域中的截断意味着我们假设：对于每个 $x$，都存在某个 $a:A(x)$ 使得 $P(x,a)$ 成立，但这些值并未以任何已知方式被选定或指定。
值域中的截断则意味着我们得出结论：存在某个函数 $g$，但这个函数并未以任何已知方式被确定或指定。

事实上，根据 \cref{thm:ttac}，这个公理也可以用一个更简单的形式来表达。

\begin{lem}\label{thm:ac-epis-split}
  选择公理~\eqref{eq:ac} 等价于下述陈述：对于任意集合 $X$ 以及任意的 $Y:X\to\type$，如果每个 $Y(x)$ 都是集合，则有
  \begin{equation}
    \Parens{\prd{x:X} \Brck{Y(x)}}
    \to
    \Brck{\prd{x:X} Y(x)}.\label{eq:epis-split}
  \end{equation}
\end{lem}

这对应于经典\index{mathematics!classical}选择公理的一个众所周知的等价形式，即``一族非空集合的 Cartesian 积是非空的''。

\begin{proof}
  由 \cref{thm:ttac}，~\eqref{eq:ac} 的值域等价于
  \[\Brck{\prd{x:X} \sm{a:A(x)} P(x,a)}.\]
  因此，~\eqref{eq:ac} 等价于将~\eqref{eq:epis-split} 应用于以下情形：\narrowequation{Y(x) \defeq \sm{a:A(x)} P(x,a).}
  （根据 \cref{thm:isset-prod,thm:prop-set}，这是一个集合。）
  反过来，~\eqref{eq:epis-split} 等价于将~\eqref{eq:ac} 应用于 $A(x)\defeq Y(x)$ 和 $P(x,a)\defeq\unit$ 的情形。
  这样，二者逻辑等价。
  由于两者都是命题，根据 \cref{lem:equiv-iff-hprop}，它们作为类型也是等价的。
\end{proof}

与 \LEM{} 类似，等价形式~\eqref{eq:ac} 和~\eqref{eq:epis-split} 并非我们基本类型论的推论，但可以相容地将其作为公理来假设。

\begin{rmk}
  不难证明，~\eqref{eq:epis-split} 的右边总是蕴含左边。
  由于两边都是命题，根据 \cref{lem:equiv-iff-hprop}，选择公理也就等价于要求一个等价
  \[ \eqv{\Parens{\prd{x:X} \Brck{Y(x)}}}{\Brck{\prd{x:X} Y(x)}} \]
  这揭示了一个常见的误区：虽然依值函数类型保持命题性（\cref{thm:isprop-forall}），但它们并不与截断交换：$\brck{\prd{x:A} P(x)}$ 通常并不等价于 $\prd{x:A} \brck{P(x)}$。
  如果我们假设选择公理，它告诉我们这对\emph{集合}是成立的；而正如我们下面将会看到的，在一般情况下它并不成立。
\end{rmk}

选择公理中对类型必须是集合的限制，可以在一定程度上放宽。
例如，我们可以允许~\eqref{eq:ac} 中的 $A$ 和 $P$，或者~\eqref{eq:epis-split} 中的 $Y$，是任意的类型族；这样会得到一个看似更强但仍然相容的陈述。
我们也可以用将在 \cref{cha:hlevels} 中考虑的更一般的 $n$-截断来替代命题截断，从而得到一族公理 $\choice n$，它们在~\eqref{eq:ac} 与~\cref{thm:ttac} 之间插值；前者简称为 \choice{}（为强调起见也可记为 $\choice{-1}$），后者我们将称为 $\choice\infty$。
亦可参见 \cref{ex:acnm,ex:acconn}。
但要注意，我们不能放宽 $X$ 必须是集合这一要求。

\begin{lem}\label{thm:no-higher-ac}
  存在一个类型 $X$ 和一个类型族 $Y:X\to \type$，使得每个 $Y(x)$ 都是集合，但条件~\eqref{eq:epis-split} 为假。
\end{lem}

\begin{proof}
  定义 $X\defeq \sm{A:\type} \brck{\bool = A}$，并设 $x_0 \defeq (\bool, \bproj{\refl{\bool}}) : X$。
  根据 $\Sigma$ 类型中道路的刻画、$\brck{A=\bool}$ 是命题这一事实以及泛等性，对任意 $(A,p),(B,q):X$，我们有 $\eqv{(\id[X]{(A,p)}{(B,q)})}{(\eqv AB)}$。
  特别地，$\eqv{(\id[X]{x_0}{x_0})}{(\eqv \bool\bool)}$，因此正如 \cref{thm:type-is-not-a-set} 中所述，$X$ 不是一个集合。

  另一方面，如果 $(A,p):X$，那么 $A$ 是一个集合。这一结论可以通过对 $p:\brck{\bool=A}$ 施行截断归纳，并利用 $\bool$ 是集合的事实得出。
  由于只要 $A$ 和 $B$ 是集合，$\eqv A B$ 就也是集合，因此对任意 $x_1,x_2:X$，$\id[X]{x_1}{x_2}$ 都是集合，也就是说 $X$ 是一个 1-类型。
  特别地，如果我们通过 $Y(x) \defeq (x_0=x)$ 来定义 $Y:X\to\UU$，那么每个 $Y(x)$ 都是集合。

  现在根据定义，对于任意 $(A,p):X$，我们有 $\brck{\bool=A}$，从而有 $\brck{x_0 = (A,p)}$。
  因此，我们得到 $\prd{x:X} \brck{Y(x)}$。
  如果条件~\eqref{eq:epis-split} 对这个 $X$ 和 $Y$ 成立，那么我们将有 $\brck{\prd{x:X} Y(x)}$。
  由于我们正试图推导出一个矛盾（$\emptyt$），而矛盾本身是一个命题，因此我们可以假设 $\prd{x:X} Y(x)$，即 $\prd{x:X} (x_0=x)$。
  但这意味着 $X$ 是一个命题，因而也是一个集合，从而产生了矛盾。
\end{proof}

\index{denial|)}%
\index{axiom!of choice|)}%

\section{唯一选择原理}
\label{sec:unique-choice}

\index{unique!choice|(defstyle}%
\indexsee{axiom!of choice!unique}{unique choice}%

以下观察虽然平凡，但非常有用。

\begin{lem}\label{thm:prop-equiv-trunc}
  如果 $P$ 是一个命题，那么 $\eqv P {\brck P}$。
\end{lem}

\begin{proof}
  当然，由定义我们有 $P\to \brck{P}$。
  并且由于 $P$ 是一个命题，将 $\brck P$ 的泛性质应用于 $\idfunc[P] :P\to P$，便得到 $\brck P \to P$。
  根据 \cref{lem:equiv-iff-hprop}，这两个函数互为拟逆。
\end{proof}

其重要推论之一如下。

\begin{cor}[唯一选择原理]\label{cor:UC}
  假设给定一个类型族 $P:A\to \type$，满足
  \begin{enumerate}
  \item 对每个 $x$，类型 $P(x)$ 都是一个命题，并且
  \item 对每个 $x$，我们有 $\brck {P(x)}$。
  \end{enumerate}
  那么我们可以得到 $\prd{x:A} P(x)$。
\end{cor}

\begin{proof}
  由上述两条假设及前一引理立即可得。
\end{proof}

这条推论还体现了一种非常有用的推理技巧。
具体来说，假设我们知道 $\brck A$，并且希望利用这一点来构造另一个类型 $B$ 的元素。
我们自然希望在构造 $B$ 的元素时能使用 $A$ 的元素，但这仅在 $B$ 是命题时才被允许，这样我们才能应用命题截断 $\brck A$ 的归纳原理；在一般情形下，我们最多只能期望证明 $\brck B$。
%
另一种做法是，我们可以为 $B$ 补充额外的数据，来\emph{唯一地}刻画我们希望构造的对象。
确切地说，我们定义一个谓词 $Q:B\to\type$，使得 $\sm{x:B} Q(x)$ 是一个命题。
这样，从 $A$ 的一个元素可以构造出满足 $Q(b)$ 的元素 $b:B$，因此从 $\brck A$ 出发可以构造 $\brck{\sm{x:B} Q(x)}$；又因为 $\brck{\sm{x:B} Q(x)}$ 等价于 $\sm{x:B} Q(x)$，即可从中投影出 $B$ 的一个元素。
这方面的例子可参见 \cref{ex:decidable-choice}。

集合论数学中也会出现类似的问题，尽管其表现方式略有不同。
如果我们试图定义一个函数 $f: A \to B$，并且对于给定的 $a : A$，我们只能证明存在某个 $b : B$，那么我们的工作尚未完成，因为我们需要具体指定一个 $B$ 的元素，而不仅仅是证明其存在性。
当然，一种方案是像我们在类型论中所做的那样，细化论证以证明 $b:B$ 的唯一存在性。
但在集合论中，通过应用选择公理——由它替我们选取所需的元素——往往可以更简单地回避这一问题。
然而在同伦类型论中，且不论是否有意避免选择公理，可用的选择公理形式本身适用性就更低，因为它们要求选择的作用域是一个\emph{集合}。
因此，如果 $A$ 不是一个集合（例如可能是宇宙 $\UU$），就不存在一致的选择公理形式，能让我们在定义 $f(a)$ 时，为每个 $a : A$ 简单地选取一个 $B$ 的元素来使用。

\index{unique!choice|)}%

\section{命题何时被截断？}
\label{subsec:when-trunc}

\index{logic!of mere propositions|(}%
\index{mere proposition|(}%
\index{logic!truncated}%

乍看之下，截断版本的 $+$ 和 $\Sigma$ 似乎比未截断的版本更接近``或''与``存在''在非形式数学中的含义。
当然，它们更接近作为形式集合论基础的一阶逻辑中\index{first-order!logic}``或''与``存在''的\emph{精确}含义，因为后者并不试图保留命题为真的任何见证。
然而，一个可能令人惊讶的事实是：非形式数学的实践往往更准确地由未截断的形式来描述。

\index{prime number}%
例如，考虑这样一个陈述：``每个素数要么是 $2$，要么是奇数''。
做研究的数学家在使用这一事实时不会感到任何不妥：他们不仅用它来证明关于素数的\emph{定理}，还用它来对素数进行\emph{构造}——在 $2$ 的情形做一件事，在奇素数的情形做另一件事。
构造的最终结果不仅仅是某个陈述的真值，而是一个可能依赖于该素数奇偶性的数据。
因此，从类型论的视角来看，这类构造很自然地要用余积类型``$(p=2)+(p\text{ is odd})$''的归纳原理来表述，而不是它的命题截断。

诚然，这不是一个最理想的例子，因为``$p=2$''与``$p$ 是奇数''是互斥的，因此 $(p=2)+(p\text{ is odd})$ 实际上已经是一个命题，从而等价于它的截断（参见 \cref{ex:disjoint-or}）。
更有说服力的例子来自存在量词。
常见的做法是：先证明一个形如``存在 $x$ 使得……''的定理，而后在后续讨论中提及``定理 Y 中构造的那个 $x$''（注意定冠词的使用）。
此外，在推导这个 $x$ 的更多性质时，可能会用到诸如``根据定理 Y 证明中对 $x$ 的构造''这样的短语。

一个非常常见的例子是``$A$ 同构于 $B$''。严格来说，这只是说 $A$ 与 $B$ 之间存在\emph{某个}同构。但在证明此类命题时，几乎总是如此：要么给出一个具体的同构，要么证明某个已知的映射是同构，而后续推理往往取决于具体给出的是哪个同构。

受过集合论训练的数学家常常会为这种``滥用语言''而感到一丝内疚。\index{abuse!of language}
我们可能会试图为此道歉，将它们从最终稿中剔除，或者用``典范的''这类模糊词语来含糊其辞。
问题还因以下事实而加剧：在形式化的集合论中，技术上根本没有办法``构造''对象——我们只能证明具有某些性质的对象存在。
因此，类型论中的\emph{非截断逻辑}捕捉到了非形式数学中一些被集合论重构所掩盖的常见做法。
（这与泛等公理佐证``将同构对象视为相同''这一常见做法——尽管它在形式上缺乏依据——是类似的。）

另一方面，\emph{截断逻辑}有时又是不可或缺的。
我们在 \LEM{} 和 \choice{} 的陈述中已经看到了这一点；本书后续还会出现其他例子。
于是，我们面临这样一个问题：在书写非形式化的类型论时，当我们使用``或''和``存在''（以及``存在一个''``我们有''这类常见同义词）这些词时，究竟应当意指什么？

达成普遍共识或许不太可能。
也许根据所从事的数学门类不同，某一种约定可能更为适用——又或许，约定如何选择根本无关紧要。
在这种情况下，在数学论文开头加一段说明，告知读者所采用的语言约定，可能就足够了。
然而，即便选定了一个总体约定，另一种逻辑通常仍会至少偶尔出现，因此我们需要一种方式来指称它。
更一般地，可以考虑用其他行为类似的类型运算来替换命题截断，例如双重否定运算 $A\mapsto \neg\neg A$，或将在 \cref{cha:hlevels} 中讨论的 $n$-截断。
作为表述方式的一种尝试，在后续内容中，我们将偶尔使用\emph{副词}\index{adverb}来表示命题截断这类``模态''的应用。

例如，如果非截断逻辑是默认约定，我们可以用副词 \define{仅仅}
\indexdef{merely}%
来表示命题截断。
于是，短语

\begin{center}
  ``仅仅存在一个 $x:A$ 使得 $P(x)$''
\end{center}

表示类型 $\brck{\sm{x:A} P(x)}$。
类似地，我们将称类型 $A$ 是\define{仅知有实例}的，
\indexdef{merely!inhabited}%
\indexdef{inhabited type!merely}%
意思是它的命题截断 $\brck A$ 是有实例的（即我们有它的一个未命名元素）。
请注意，这是对副词``merely''在非形式数学英语中用法的\emph{定义}\index{definition!of adverbs}，
正如我们为``群''和``环''等名词\index{noun}，以及``正则的''和``正规的''等形容词\index{adjective}赋予精确数学含义一样。
我们并非声称``merely''在词典中的定义就是指命题截断；选择这个词只是意在提醒数学家读者：一个命题所包含的信息``仅仅''是一个真值，再无其他。

另一方面，如果截断逻辑是当前的默认约定，我们可以用\define{纯粹地}
\indexdef{purely}%
或\define{构造性地}这样的副词来标示截断的缺失，从而

\begin{center}
``纯粹地存在一个 $x:A$ 使得 $P(x)$''
\end{center}

将表示类型 $\sm{x:A} P(x)$。
即使无截断已是默认约定，我们也可以用``purely''或``actually''来强调这一点。

在本书中，出于若干原因，我们将继续以非截断逻辑作为默认约定。

\begin{enumerate}[label=(\arabic*)]
\item 我们想鼓励初学者去尝试使用它，而不是因为更熟悉就固守截断逻辑。
\item 在类型论中将截断逻辑作为默认，会面临与集合论基础同类的``语言滥用''\index{abuse!of language}问题，而未截断逻辑则能避免这些问题。
  举例来说，我们将 ``$\eqv A B$'' 定义为 $A$ 与 $B$ 之间等价的类型，而非其命题截断，这意味着证明形如 ``$\eqv A B$'' 的定理，字面上就是构造一个具体的此类等价。
  这个具体的等价之后可以引用。
\item 我们想强调，``命题''这一概念并非类型论的基本组成部分。
  正如我们将在 \cref{cha:hlevels} 中看到的，命题仅仅是无限阶梯中的第二级，此外还存在许多完全不在这个阶梯上的其他模态。
\item 许多在经典意义下是命题的陈述，在同伦类型论中不再如此。
  当然，其中最重要的便是相等。
\item 另一方面，同伦类型论最有趣的观察之一是，有数量惊人的类型\emph{自动}成为命题，或者稍加修改即可如此，而无需任何截断。
  （参见 \cref{thm:isprop-isprop,cha:equivalences,cha:hlevels,cha:category-theory,cha:set-math}。）
  因此，尽管这些类型除真值外不含任何数据，我们仍然可以用它们来构造未截断的对象，因为无需使用命题截断的归纳原理。
  如果默认对所有陈述都应用命题截断，这一有用事实的表达会变得更为笨拙。
\item 最后，截断对于本书中的大部分数学来说用处不大，因此当它们出现时明确标注出来更为简便。
\end{enumerate}

\index{mere proposition|)}%
\index{logic!of mere propositions|)}%

\section{可缩性}
\label{sec:contractibility}

\index{type!contractible|(defstyle}%
\index{contractible!type|(defstyle}%

在 \cref{thm:inhabprop-eqvunit} 中我们观察到，一个有实例的命题必然等价于 $\unit$，
\index{type!unit}%
反过来也成立，这不难看出。
具有这一性质的类型称为\emph{可缩的}。
可缩性的另一个等价定义，有时也较为方便，如下所述。

\begin{defn}\label{defn:contractible}
  一个类型 $A$ 是\define{可缩的（contractible）}，
  或称\define{单点集}，
  \indexdef{type!singleton}%
  \indexsee{singleton type}{type, singleton}%
  如果存在 $a:A$，称为\define{收缩中心}，
  \indexdef{center!of contraction}%
  使得对所有 $x:A$ 都有 $a=x$。
  我们将指定的道路 $a=x$ 记作 $\contr_x$。
\end{defn}

换言之，类型 $\iscontr(A)$ 定义为
\[ \iscontr(A) \defeq \sm{a:A} \prd{x:A}(a=x). \]
注意，在通常的命题即类型解读下，我们可以将 $\iscontr(A)$ 读作``$A$ 恰好包含一个元素''，或更精确地说，``$A$ 包含一个元素，且 $A$ 的每个元素都等于该元素''。

\begin{rmk}
  我们还可以从拓扑角度将 $\iscontr(A)$ 读作``存在点 $a:A$，使得对所有 $x:A$，存在一条从 $a$ 到 $x$ 的道路''。
  注意，对熟悉经典数学的人而言，这听起来更像是\emph{连通性}而非可缩性的定义。
  \index{continuity of functions in type theory@``continuity'' of functions in type theory}%
  \index{functoriality of functions in type theory@``functoriality'' of functions in type theory}%
  关键在于，此句中``存在''的含义是连续的/自然的。

  表达连通性的更好方式是 $\sm{a:A}\prd{x:A} \brck{a=x}$。
  如果假定 $A$ 是带点的，这确实是正确的——参见 \cref{thm:connected-pointed} 后的备注——但一般而言，一个类型可以不带基点而连通。
  在 \cref{sec:connectivity} 中，我们将把连通性定义为一般 $n$-连通性概念在 $n=0$ 时的特例，并且在 \cref{ex:connectivity-inductively} 中，请读者证明该定义等价于同时拥有 $\brck{A}$ 与 $\prd{x,y:A} \brck{x=y}$。
\end{rmk}

\begin{lem}\label{thm:contr-unit}
  对于类型 $A$，以下逻辑等价。
  \begin{enumerate}
  \item $A$ 是按 \cref{defn:contractible} 定义的可缩的。\label{item:contr}
  \item $A$ 是一个命题，并且存在点 $a:A$。\label{item:contr-inhabited-prop}
  \item $A$ 等价于 \unit。\label{item:contr-eqv-unit}
  \end{enumerate}
\end{lem}

\begin{proof}
  如果 $A$ 可缩，那么它显然有一个点 $a:A$（即收缩中心），并且对于任意 $x,y:A$ 都有 $x=a=y$；因此 $A$ 是一个命题。
  反过来，如果我们有一个点 $a:A$ 并且 $A$ 是一个命题，那么对于任意 $x:A$ 都有 $x=a$；因此 $A$ 可缩。
  而我们已在 \cref{thm:inhabprop-eqvunit} 中证明了~\ref{item:contr-inhabited-prop}$\Rightarrow$\ref{item:contr-eqv-unit}，其逆命题则可由单位类型~\unit 显然满足性质~\ref{item:contr-inhabited-prop} 得出。
\end{proof}

\begin{lem}\label{thm:isprop-iscontr}
  对于任何类型 $A$，类型 $\iscontr(A)$ 是一个命题。
\end{lem}

\begin{proof}
  假设给定 $c,c':\iscontr(A)$。
  我们可设 $c\jdeq(a,p)$ 与 $c'\jdeq(a',p')$，其中 $a,a':A$，且有 $p:\prd{x:A} (a=x)$ 与 $p':\prd{x:A} (a'=x)$。
  根据 $\Sigma$-类型中道路的特征，要证明 $c=c'$，只需给出 $q:a=a'$ 使得 $\trans{q}{p}=p'$。
  %
  我们取 $q\defeq p(a')$。
  现在，由于 $A$ 可缩（根据 $c$ 或 $c'$），由 \cref{thm:contr-unit} 可知它是一个命题。
  因此由 \cref{thm:prop-set,thm:isprop-forall}，$\prd{x:A}(a'=x)$ 也是命题；于是 $\trans{q}{p}=p'$ 自动成立。
\end{proof}

\begin{cor}\label{thm:contr-contr}
  如果 $A$ 可缩，那么 $\iscontr(A)$ 也可缩。
\end{cor}

\begin{proof}
  由 \cref{thm:isprop-iscontr} 与 \cref{thm:contr-unit}\ref{item:contr-inhabited-prop} 即得。
\end{proof}

如同命题一样，可缩类型也被许多类型构造子所保持。
例如，我们有：

\begin{lem}\label{thm:contr-forall}
  如果 $P:A\to\type$ 是一个类型族，且每个 $P(a)$ 都可缩，那么 $\prd{x:A} P(x)$ 也可缩。
\end{lem}

\begin{proof}
  根据 \cref{thm:isprop-forall}，$\prd{x:A} P(x)$ 是一个命题，因为每个 $P(x)$ 都是命题。
  但它也有一个元素，即把每个 $x:A$ 映到 $P(x)$ 的收缩中心的函数。
  于是由 \cref{thm:contr-unit}\ref{item:contr-inhabited-prop} 可知，$\prd{x:A} P(x)$ 可缩。
\end{proof}

\index{function extensionality}%
（事实上，\cref{thm:contr-forall} 的陈述等价于函数外延性公理。
参见~\cref{sec:univalence-implies-funext}。）

显然，如果 $A$ 等价于 $B$ 且 $A$ 可缩，那么 $B$ 也可缩。
更一般地，只需 $B$ 是 $A$ 的一个\emph{收缩核（retract）}即可。
按定义，一个\define{收缩}
\indexdef{retraction}%
\indexdef{function!retraction}%
是一个函数 $r : A \to B$，使得存在一个函数 $s : B \to A$（称为它的\define{截面}，
\indexdef{section}%
\indexdef{function!section}%）
以及一个同伦 $\epsilon:\prd{y:B} (r(s(y))=y)$；此时我们称 $B$ 是 $A$ 的一个\define{收缩核}%
\indexdef{retract!of a type}。

\begin{lem}\label{thm:retract-contr}
  如果 $B$ 是 $A$ 的一个收缩核，并且 $A$ 是可缩的，那么 $B$ 也是可缩的。
\end{lem}

\begin{proof}
  设 $a_0 : A$ 为收缩中心。
  我们断言 $b_0 \defeq r(a_0) : B$ 是 $B$ 的一个收缩中心。
  令 $b : B$；我们需要一条道路 $b = b_0$。
  但我们有 $\epsilon_b : r(s(b)) = b$ 与 $\contr_{s(b)} : s(b) = a_0$，因此通过复合可得
  \[ \opp{\epsilon_b} \ct \ap{r}{\contr_{s(b)}} : b = r(a_0) \jdeq b_0. \qedhere\]
\end{proof}

可缩类型乍看可能并不十分有趣，因为它们都等价于 \unit。
这一概念之所以有用，原因之一在于，有时一组各自非平凡的数据，整体上却会形成一个可缩的类型。
一个重要的例子就是带一个自由端点的道路空间。
正如我们将在 \cref{sec:identity-systems} 中看到的，这个事实本质上封装了相等类型中基于一点的道路归纳原理。

\begin{lem}\label{thm:contr-paths}
  对任意 $A$ 与任意 $a:A$，类型 $\sm{x:A} (a=x)$ 是可缩的。
\end{lem}

\begin{proof}
  我们选择点 $(a,\refl a)$ 作为中心。
  现在假设 $(x,p):\sm{x:A}(a=x)$；我们必须证明 $(a,\refl a) = (x,p)$。
  根据 $\Sigma$-类型中道路的特征，只需给出 $q:a=x$，使得 $\trans{q}{\refl a} = p$。
  但我们可取 $q\defeq p$，此时 $\trans{q}{\refl a} = p$ 由道路类型中输运的特征可得。
\end{proof}

出现这种情况时，我们便可在等价意义下简化复杂的构造，所依据的非正式原则是：可缩的数据可以自由忽略。
这一原则体现为许多引理，其中大部分留给读者；下面是一例。

\begin{lem}\label{thm:omit-contr}
  设 $P:A\to\type$ 是一个类型族。
  \begin{enumerate}
  \item 如果每个 $P(x)$ 都是可缩的，那么 $\sm{x:A} P(x)$ 等价于 $A$。\label{item:omitcontr1}
  \item 如果 $A$ 是可缩的且收缩中心为 $a$，那么 $\sm{x:A} P(x)$ 等价于 $P(a)$。\label{item:omitcontr2}
  \end{enumerate}
\end{lem}

\begin{proof}
  对于情况~\ref{item:omitcontr1}，我们证明 $\proj1:\sm{x:A} P(x) \to A$ 是一个等价。
  取其拟逆为 $g(x)\defeq (x,c_x)$，其中 $c_x$ 是 $P(x)$ 的收缩中心。
  复合 $\proj1 \circ g$ 显然是 $\idfunc[A]$，而反向复合利用每个 $P(x)$ 的收缩性，同伦于恒等映射。

  我们将情况~\ref{item:omitcontr2} 的证明留给读者（参见 \cref{ex:omit-contr2}）。
\end{proof}

可缩类型之所以有趣的另一个原因在于，它们将 \cref{sec:basics-sets} 中提到的 $n$-类型的阶梯向下又延伸了一级。

\begin{lem}\label{thm:prop-minusonetype}
  一个类型 $A$ 是命题，当且仅当对所有 $x,y:A$，类型 $\id[A]xy$ 都是可缩的。
\end{lem}

\begin{proof}
  对于“如果”方向，我们只需注意到任何可缩类型都是有实例的。
  对于“仅当”方向，我们在 \cref{subsec:hprops} 中已观察到每个命题都是集合，因而每个类型 $\id[A]xy$ 都是命题。
  但它也是有实例的（因为 $A$ 是命题），于是由 \cref{thm:contr-unit}\ref{item:contr-inhabited-prop} 可知它是可缩的。
\end{proof}

因此，可缩类型也可以称为 \define{$(-2)$-类型（$(-2)$-types）}。
它们是 $n$-类型阶梯的最底层，也将是 \cref{cha:hlevels} 中递归定义 $n$-类型时的基础情形。

\index{type!contractible|)}%
\index{contractible!type|)}%

\sectionNotes

在类型论中可以定义集合、命题和可缩类型，且所有高阶同伦都如 \cref{sec:basics-sets,subsec:hprops,sec:contractibility} 中那样自动得到处理——这一事实最早由 Voevodsky 注意到。
事实上，他通过归纳定义了整个 $n$-类型的层级，正如我们将在 \cref{cha:hlevels} 中所做的那样。\index{n-type@$n$-type!definable in type theory}%

\cref{thm:not-dneg,thm:not-lem} 本质上依赖于 Hedberg 的一个经典定理，
\index{Hedberg's theorem}%
\index{theorem!Hedberg's}%
我们将在 \cref{sec:hedberg} 中加以证明。
“命题即类型”形式的 \LEM{} 与泛等性相矛盾，这一结论是由 Martín Escardó（马丁·埃斯卡多）在 \Agda 邮件列表上指出的。
而我们给出的 \cref{thm:not-dneg} 的证明则归功于 Thierry Coquand。

命题截断最初由 Constable~\cite{Con85} 于 1983 年作为“子集”类型和“商”类型的一种应用，引入到 \NuPRL 的外延类型论
\index{proof!assistant!\NuPRL}%
之中。在 \NuPRL 类型论~\cite{constable+86nuprl-book} 中，本书所称的“命题截断”当时被称为“压平”。
此后，~\cite{ab:bracket-types} 仍在外延类型论的框架下，给出了直接刻画命题截断的规则。而在同伦类型论中的内涵版本，则由 Voevodsky 利用非直谓\index{impredicative!truncation}量化构造，随后 Lumsdaine 又用高阶归纳类型加以构造（参见 \cref{sec:hittruncations}）。

\index{propositional!resizing}%
Voevodsky~\cite{Universe-poly} 提出了 \cref{subsec:prop-subsets} 中所讨论的那一类大小重整规则。
\index{axiom!of reducibility}\index{resizing}%
这些规则显然与 Russell 在他与 Whitehead 合著的《数学原理》~\cite{PM2} 中提出的那个备受争议的\emph{可化归公理（axiom of reducibility）}密切相关。\index{Russell, Bertrand}

副词“纯粹地（purely）”用于指代未截断的逻辑，这一用法是在援引编程语言中以单子模态来建模副作用（effects）的做法；参见 \cref{sec:modalities} 及 \cref{cha:hlevels} 的附注。

相对于类型论，逻辑可以有多种不同的处理方式。
例如，除了 \cref{sec:pat} 中描述的单纯的“命题即类型”逻辑，以及 \cref{subsec:logic-hprop} 描述的仅使用命题的另一种方式外，人们还可以引入一种独立的命题 ``sort''，其行为在某些方面类似于类型，但不与类型等同。
此方法被逻辑充实类型论~\cite{aczel2002collection}、对意象（topos）及相关范畴内部语言的某些表述（例如~\cite{jacobs1999categorical,elephant}）以及证明助理 \Coq\index{proof!assistant!Coq@\textsc{Coq}} 所采用。
这种处理方法更一般，但功能也较弱。
举例来说，唯一选择原理（\cref{sec:unique-choice}）在 \Coq~\cite{Spiwack} 中所谓的 setoid 构成的范畴里、在逻辑充实类型论~\cite{aczel2002collection} 中、以及在最小类型论~\cite{maietti2005toward} 中都是不成立的。\index{setoid}
因此，泛等公理使我们的类型论更接近意象的内部逻辑；另见 \cref{cha:set-math}。\index{topos}

Martin-L\"of~\cite{martin2006100} 对选择公理的历史进行了讨论。当然，构造性数学和直觉主义数学有着悠久而复杂的历史，我们在此不再深究；例如可参阅~\cite{TroelstraI,TroelstraII}。

\sectionExercises

\begin{ex}\label{ex:equiv-functor-set}
  证明若 $\eqv A B$ 且 $A$ 是一个集合，则 $B$ 亦然。
\end{ex}

\begin{ex}\label{ex:isset-coprod}
  证明若 $A$ 和 $B$ 都是集合，则 $A+B$ 也是集合。
\end{ex}

\begin{ex}\label{ex:isset-sigma}
  证明若 $A$ 是一个集合，且 $B:A\to \type$ 是一个类型族，使得对所有 $x:A$，$B(x)$ 都是集合，则 $\sm{x:A} B(x)$ 也是集合。
\end{ex}

\begin{ex}\label{ex:prop-endocontr}
  证明 $A$ 是一个命题当且仅当 $A\to A$ 是可缩的。
\end{ex}

\begin{ex}\label{ex:prop-inhabcontr}
  证明 $\eqv{\isprop(A)}{(A\to\iscontr(A))}$。
\end{ex}

\begin{ex}\label{ex:lem-mereprop}
  证明若 $A$ 是一个命题，则 $A+(\neg A)$ 也是命题。
  因此，在~\eqref{eq:lem} 处无需插入命题截断。
\end{ex}

\begin{ex}\label{ex:disjoint-or}
  更一般地，证明若 $A$ 和 $B$ 都是命题且 $\neg(A\times B)$，则 $A+B$ 也是命题。
\end{ex}

% \begin{ex}\label{ex:hprop-iff-equiv}
%   Show that if $A$ and $B$ are mere propositions such that $A\to B$ and $B\to A$, then $\eqv A B$.
% \end{ex}

% \begin{ex}\label{ex:isprop-isprop}
%   Show that for any type $A$, the types $\isprop(A)$ and $\isset(A)$ are mere propositions.
% \end{ex}

% \begin{ex}\label{ex:prop-eqvtrunc}
%   Show that if $A$ is already a mere proposition, then $\eqv A{\brck{A}}$.
% \end{ex}

\begin{ex}\label{ex:brck-qinv}
  假设某个类型 $\isequiv(f)$ 满足\cref{sec:basics-equivalences}的条件~\ref{item:be1}--\ref{item:be3}，证明类型 $\brck{\qinv(f)}$ 满足同样的条件，并且等价于 $\isequiv(f)$。
\end{ex}

\begin{ex}\label{ex:lem-impl-prop-equiv-bool}
  证明：若排中律成立，则类型$\prop \defeq \sm{A:\type} \isprop(A)$等价于 \bool。
\end{ex}

\begin{ex}\label{ex:lem-impred}
  证明：若 $\UU_{i+1}$ 满足排中律，则典范包含$\prop_{\UU_i} \to \prop_{\UU_{i+1}}$是一个等价。
\end{ex}

\begin{ex}\label{ex:not-brck-A-impl-A}
  证明：并非对所有的 $A:\type$ 都有 $\brck{A} \to A$。
  （不过，对于某些特定的类型可以有 $\brck{A}\to A$。
  \cref{ex:brck-qinv}蕴含 $\qinv(f)$ 就是这样的类型。）
\end{ex}

\begin{ex}\label{ex:lem-impl-simple-ac}
  \index{axiom!of choice}%
  证明：若排中律成立，则对所有的 $A:\type$，有 $\bbrck{(\brck A \to A)}$。
  （这个性质是选择公理的一种极简形式，在没有排中律的情况下可能不成立；参见~\cite{krausgeneralizations}。）
\end{ex}

\begin{ex}\label{ex:naive-lem-impl-ac}
  我们在\cref{thm:not-lem}中证明了，下述朴素形式的排中律与泛等性不相容：
  \[ \prd{A:\type} (A+(\neg A)) \]
  在没有泛等性的情况下，这条公理是一致的。
  然而，证明它蕴含选择公理~\eqref{eq:ac}。
\end{ex}

\begin{ex}\label{ex:lem-brck}
  证明：假设排中律成立，双重否定 $\neg \neg A$ 与命题截断 $\brck A$ 具有相同的递归原理，但其计算规则是命题性的而非判断性的。
  换言之，证明：假设排中律成立，若 $B$ 是一个命题且有 $f:A\to B$，则存在导出的 $g:\neg\neg A \to B$ 使得对所有 $a:A$ 都有 $g(\bproj a) = f(a)$。
  由此推出（假设排中律成立）$\eqv{\neg\neg A}{\brck{A}}$。
  因此，在排中律下，命题截断可以被定义出来，而不必作为一个单独的类型构造子。
\end{ex}

\begin{ex}\label{ex:impred-brck}\ 
  \index{propositional!resizing}%
  \begin{enumerate}
  \item 证明：对任何 $A : \UU$，类型
    \[\prd{P:\prop_{\UU}} \Parens{(A\to P)\to P}\]
    与 $\brck A$ 具有相同的递归原理（至少相对于 $\UU$ 中的命题而言），\emph{并}具有相同的判断性计算规则。
  \item 上一部分考虑的类型并不位于同一个宇宙 $\UU$ 中，而且它的递归原理只适用于 $\UU$ 中的命题。
    证明：如果假设命题大小重整，那么可以定义一个确实位于同一宇宙 $\UU$ 中的类型，它满足与 $\brck A$ 相同的递归原理，尽管其计算规则只能是命题性的。
    因此，在这种情况下我们同样可以定义命题截断。
  \end{enumerate}
\end{ex}

\begin{ex}\label{ex:lem-impl-dn-commutes}
  假设排中律成立，证明双重否定与集合上命题的全称量化可交换。
  也就是说，证明：若 $X$ 是一个集合且每个 $Y(x)$ 都是命题，则排中律蕴含
  \begin{equation}
    \eqv{\Parens{\prd{x:X} \neg\neg Y(x)}}{\Parens{\neg\neg \prd{x:X} Y(x)}}.\label{eq:dnshift}
  \end{equation}
  注意，如果我们改为假设每个 $Y(x)$ 都是集合，那么~\eqref{eq:dnshift}将等价于选择公理~\eqref{eq:epis-split}。
\end{ex}

\begin{ex}\label{ex:prop-trunc-ind}
  \index{induction principle!for truncation}%
  证明\cref{subsec:prop-trunc}中给出的命题截断规则足以推出下述归纳原理：对于任意类型族 $B:\brck A \to \type$，若每个 $B(x)$ 都是命题，且对每个 $a:A$ 都有 $B(\bproj a)$，则对每个 $x:\brck A$ 都有 $B(x)$。
\end{ex}

\begin{ex}\label{ex:lem-ldn}
  证明排中律~\eqref{eq:lem}与双重否定律~\eqref{eq:ldn}是逻辑等价的。
\end{ex}

\begin{ex}\label{ex:decidable-choice}
  设 $P:\nat\to\type$ 是一个可判定的命题族（参见\cref{defn:decidable-equality}\ref{item:decidable-equality2}）。
  证明：
  \[ \Brck{\sm{n:\nat} P(n)} \;\to\; \sm{n:\nat}P(n).\]
\end{ex}

\begin{ex}\label{ex:omit-contr2}
  证明\cref{thm:omit-contr}\ref{item:omitcontr2}：如果 $A$ 是可缩的，且收缩中心为 $a$，那么 $\sm{x:A} P(x)$ 等价于 $P(a)$。
\end{ex}

\begin{ex}\label{ex:isprop-equiv-equiv-bracket}
  证明$\isprop(P) \simeq (P \simeq \brck P)$。
\end{ex}

\begin{ex}\label{ex:finite-choice}
  如同在经典集合论中一样，选择公理的有限版本是一个定理。证明当 $X$ 是有限类型 $\Fin(n)$（定义见\cref{ex:fin}）时，选择公理\eqref{eq:ac}成立。
\end{ex}

\begin{ex}\label{ex:decidable-choice-strong}
  证明：如果 $P:\nat\to\type$ 是任意可判定族，则\cref{ex:decidable-choice}的结论成立。
\end{ex}

\begin{ex}\label{ex:n-set}
  通过先证明对所有 $m,n$，$\code(m,n)$ 都是命题，来简化\cref{thm:path-nat}的证明。
\end{ex}

% Local Variables:
% TeX-master: "hott-online"
% End: